% SPDX-FileCopyrightText: 2026 Bart van de Lint
% SPDX-License-Identifier: AGPL-3.0-or-later
\section{Details of the discussion}
\label{sec:disc_details}

\subsection{Ship kites at sea}
\label{sec:disc_details_kites}

SkySails reported more than \num{1200} automated kite flights and more than \num{1800} hours of air time over all its kites, from 6 to \SI{400}{m^2}, with single flights of up to 120 hours on a ship \citep{skysails2018}. The kites were ram-air wings, launched and recovered by a telescopic crane that held them at the leading edge and inflated them. The project record reports a saving of 5\,\% on the average route mix \citep{wintecc2009_record}, with no kite size. Our model gives \tbd{1.7}\,\% for a kite of \SI{160}{m^2}, the size the Beluga SkySails was launched with \citep{wiwo2008}, on a ship of the same size (\cref{sec:checks}), and a kite of up to \SI{320}{m^2} later flew on that ship \citep{wrage2018}.

The trial kites of 250 and \SI{500}{m^2} of the French maker Airseas \citep{airseas_cp2022} were double-skinned ram-air wings too \citep{noaa2009,airseas_seawing_web}. Leading-edge inflatable kites, whose shape comes from air-filled tubes along the front edge and across the canopy \citep{poland2025}, have flown for power at 60 to \SI{80}{m^2} flat as a hybrid with a fibre-glass skeleton \citep{kitepower2024}, and no source reports one of ship-kite size.

\subsection{Series prices}
\label{sec:disc_details_prices}

\textbf{Check against kite power.} In a ground-generation kite power system, where the kite pulls its tether off a drum that drives a generator on land, the same tether and control unit and a canopy at today's price give a cost of electricity of about 250~EUR/MWh, dominated by the canopy replaced every \num{2000} flight hours. Published estimates for such systems are 18 to 150~EUR/MWh \citep{joshi2025}, so the components are not priced too low for what they are.

\textbf{Business case of a kite maker.} A maker that spends about USD~100 million on development and sells 200 systems of \SI{800}{m^2} a year breaks even after about thirteen years at today's canopy price and the canopy life goal of \num{3000} flight hours, with a net present value at 10\,\% of USD~61 million. The price goal moves this to fourteen years and minus USD~18 million. Canopies lasting \num{4000} flight hours as well, when a ship buys fewer, move it to fifteen years and minus USD~47 million. At the pace of today's rotor makers it does not break even.

\subsection{Interference between deck rigs}
\label{sec:disc_details_size}

In two-dimensional potential flow, three wing sails in line lose up to 6\,\% of thrust and four in a two-by-two layout 10 to 28\,\% \citep{vanreen2026}. In the wind tunnel, two rotors three diameters apart in line lost up to 37\,\% of driving force at \SI{30}{\degree} apparent wind, gained up to 17\,\% downwind, and lost 4 to 6\,\% on average over the wind angles of a reaching ship \citep{bordogna2020}. CFD of wing sails on a bulk carrier gives 98\,\% of the isolated thrust for three sails 2.4 chords apart and 77\,\% for five sails 1.4 chords apart \citep{mboumba2022}, the chord being the depth of a wing from front to back. Our array factor gives 91 and 85\,\%, which takes back a third of the loss for sails trimmed one by one, and 87 and 78\,\% at fixed trim.

\subsection{Wind shear by region}
\label{sec:disc_details_profile}

In the well-mixed trade-wind belt our ERA5 profiles give only 1.17 times the \SI{10}{m} wind at \SI{540}{m}. North of \SI{40}{\degree N} they follow the logarithmic law with \SI{0.2}{mm} roughness, an exponent of about 0.08 between \SI{10}{m} and \SI{270}{m}, which matches the four-year mean of mast and lidar measurements \SI{85}{km} off the Dutch coast up to \SI{315}{m} better than a power law with the usual offshore exponent of 0.14 \citep{kalverla2020thesis}.

\subsection{Calm-water resistance of full hulls}
\label{sec:disc_details_hull}

\citet{holtrop1982} limit their method to hulls resembling their average ship and call its high-block adjustments tentative. \citet{ortigosa2009} tabulate a prismatic coefficient of at most 0.85 for tankers and bulk carriers, a limit we could read only second-hand. Beyond it the fitted wave term, the part of the resistance due to the waves the hull makes, grows steeply. Uncalibrated it takes 25 and 31\,\% of the resistance of the Ultramax and Kamsarmax at service speed, and one calibration factor scales every component, which is what would make their power at operating speed too low.
